% ============================================================
% IEEE CONFERENCE PAPER — IEEEtran class, two-column
% Title: Near-Miss Citation Verification in Indian Supreme Court
%        Judgments via Complexity-Adaptive RAG
% ============================================================
\documentclass[conference]{IEEEtran}

% ---------- packages ----------
\usepackage[T1]{fontenc}
\usepackage[utf8]{inputenc}
\usepackage{cite}
\usepackage{amsmath,amssymb,amsfonts}
\usepackage{graphicx}
\usepackage{textcomp}
\usepackage{xcolor}
\usepackage{booktabs}
\usepackage{multirow}
\usepackage{url}
\usepackage{microtype}
\usepackage{hyperref}
\hypersetup{colorlinks=false,pdfborder={0 0 0}}

% ---------- IEEE-required fix ----------
\def\BibTeX{{\rm B\kern-.05em{\sc i\kern-.025em b}\kern-.08em
    T\kern-.1667em\lower.7ex\hbox{E}\kern-.125emX}}

% ============================================================
\begin{document}

\title{Near-Miss Citation Verification in Indian Supreme Court Judgments\\
via Complexity-Adaptive RAG}

\author{
  \IEEEauthorblockN{Nishchal K J}
  \IEEEauthorblockA{\textit{Dept. of Computer Science \& Engineering}\\
  \textit{Ramaiah Institute of Technology}\\
  Bangalore, India\\
  nishchal@example.edu}
  \and
  \IEEEauthorblockN{Riya Acharya}
  \IEEEauthorblockA{\textit{Dept. of Computer Science \& Engineering}\\
  \textit{Ramaiah Institute of Technology}\\
  Bangalore, India\\
  riya@example.edu}
  \and
  \IEEEauthorblockN{Saksham Yadav}
  \IEEEauthorblockA{\textit{Dept. of Computer Science \& Engineering}\\
  \textit{Ramaiah Institute of Technology}\\
  Bangalore, India\\
  saksham@example.edu}
  \and
  \IEEEauthorblockN{Sangeetha V}
  \IEEEauthorblockA{\textit{Dept. of Computer Science \& Engineering}\\
  \textit{Ramaiah Institute of Technology}\\
  Bangalore, India\\
  drsangeethav@msrit.edu}
}

\maketitle

% ============================================================
\begin{abstract}
Citation errors in legal AI go beyond non-existent cases: a cited authority may exist yet fail
to support the proposition. We study this in Indian law and contribute three items.
\textbf{(i)~Benchmark.} \textbf{IndiaLegalCite} comprises 286 (v0) and 432 (v1) citation-support
items from 199 Supreme Court of India judgments (2023--2024), covering eight corruption types
including three novel \emph{near-miss} categories and an IPC/BNS statute-transition error unique
to Indian jurisdiction.
\textbf{(ii)~Baselines.} Existence checks achieve only F1\,=\,0.31 on v1; local NLI reaches
F1\,=\,0.86 but falsely flags 31\% of faithful claims and misses 48\% of over-broad holdings;
confidence-based abstention raises accuracy to 0.99 at 20\% coverage.
\textbf{(iii)~Adaptive RAG.} A hybrid complexity-scoring pipeline with sub-question decomposition,
Reciprocal Rank Fusion, and NLI+LLM verification correctly modulates retrieval depth but yields
no accuracy gain over a fixed top-5 baseline (McNemar $p=0.0625$), at 13.5$\times$ latency.
Code and benchmark manifests are released openly.
\end{abstract}

\begin{IEEEkeywords}
legal NLP, citation verification, retrieval-augmented generation, Indian Supreme Court,
near-miss errors, adaptive retrieval, evidence verification, negative results
\end{IEEEkeywords}

% ============================================================
\section{Introduction}
\label{sec:intro}

When a legal AI system cites a case, two things need to be true: the case must exist, and it must
actually support what the system claims. Most tools only check the first part. Standard RAG
pipelines retrieve passages by vector similarity and generate an answer; nothing checks whether the
cited passage genuinely supports the claim~\cite{magesh2025hallucination}.

Verma~\cite{verma2026citation} has studied this for US court opinions and found that wrong-case
substitutions are caught reliably (93--100\% recall) but wrong-pinpoint errors are frequently
missed (37--61\%). GANDR~\cite{qian2026gandr} pairs an answer-writing agent with an auditing critic
and reports 70.8\% strict accuracy and F1 = 0.84 on 185 items. Neither study covers Indian law,
its citation conventions, or the IPC/BNS statute replacement of July 2024.

\textbf{Contributions:}
\begin{enumerate}
\item \textbf{IndiaLegalCite benchmark}: two versioned datasets of citation-support items from
  Indian Supreme Court judgments, with three novel near-miss error types and an IPC/BNS
  code-transition error specific to Indian law.
\item \textbf{Baseline evaluation}: existence checks catch only fabricated citations; a local
  NLI model is viable but imprecise; confidence-based abstention substantially improves reliability.
\item \textbf{Complexity-adaptive RAG pipeline}: a continuous hybrid complexity score drives
  retrieval depth; sub-question decomposition with RRF broadens coverage; NLI+LLM verification
  closes the loop. We report an honest negative result: depth adaptation works, accuracy does not
  improve, and cost is $13.5\times$ higher.
\end{enumerate}

% ============================================================
\section{Related Work}
\label{sec:related}

Magesh et al.~\cite{magesh2025hallucination} showed that commercial legal tool hallucination-free
claims are overstated. Verma~\cite{verma2026citation} built a citation-perturbation framework
showing topical overlap and legal support are distinct model capabilities. GANDR~\cite{qian2026gandr}
achieves F1 = 0.84 on legal claim auditing with a critic-augmented pipeline.

For factual verification, FActScore~\cite{min2023factscore} popularised atomic-claim decomposition.
ALCE~\cite{gao2023alce} defines NLI-based citation precision/recall; RARR~\cite{gao2023rarr} edits
claims post-hoc when unsupported. Our verifier draws on all three, adding a \emph{quote-check}
guard against LLM-hallucinated evidence quotes.

For adaptive retrieval, Adaptive-RAG~\cite{jeong2024adaptiverag} uses a trained three-way
classifier; Self-RAG~\cite{asai2024selfrag} learns to self-critique; CRAG~\cite{yan2024crag} adds
a corrective retrieval phase. LegalBench-RAG~\cite{pipitone2024legalbench} benchmarks legal
retrieval. Our approach uses a \emph{continuous} score (four sub-dimensions, smooth $k$
interpolation) with no specialised training. InLegalNLP~\cite{paul2022inlegal} and
Legal-BERT~\cite{chalkidis2020legalbert} provide Indian and European legal representations, but
neither addresses citation-level support checking.

% ============================================================
\section{IndiaLegalCite Benchmark}
\label{sec:benchmark}

\textbf{Corpus.} We sampled 200 English Supreme Court of India judgments at random (seed 13;
100 per year) from the AWS Open Data registry~\cite{indiansc2026} and downloaded 199 successfully.
Paragraphs are split on the court's sequential numbering (84\% of judgments); blank lines for the
rest. Mean 27.2 paragraphs per judgment. The licence is listed as CC-BY-4.0 but applicability to
judgment texts is unconfirmed; we do not redistribute any judgment text.

\textbf{Items and taxonomy.} An item is a triple: a claim, a cited authority, and a binary flag
label. Each source proposition yields one faithful control and one or more corruptions, matched
by source. Table~\ref{tab:composition} shows the counts. \emph{Surface-level} error types (v0 and
v1): fabricated case, wrong paragraph, altered quote (negation/number/modal), and wrong
code (IPC/BNS statute replacement, 1 July 2024). \emph{Near-miss} types (v1 only, designed to
defeat topical-overlap acceptance): near paragraph (most similar other paragraph of same judgment),
over-broad holding (condition word deleted; 75 candidates rejected by filter), and swapped
party/outcome.

\textbf{Generation and quality.} v0 used \texttt{gpt-4o-mini}; v1 used \texttt{gpt-4.1} with
\texttt{gpt-4.1-mini} as filter. Items are rejected if the supporting quote is not verbatim in
the source, or if the filter judges a corrupted claim still supported. Labels are automatic by
construction (not lawyer-validated). One author read 30 stratified v1 items: none clearly
mislabelled, 2 (7\%) borderline. The IPC/BNS mapping table is partial and unverified.

\begin{table}[t]
\centering
\caption{IndiaLegalCite benchmark composition (dev / test, by source judgment).}
\label{tab:composition}
\renewcommand{\arraystretch}{1.1}
\small
\begin{tabular}{lrrrr}
\toprule
\textbf{Item type} & \textbf{v0 dev} & \textbf{v0 test} & \textbf{v1 dev} & \textbf{v1 test} \\
\midrule
Faithful control      & 18 & 44  & 32 & 111 \\
Fabricated case       & 15 & 35  &  8 &  42 \\
Wrong paragraph       & 15 & 35  &  8 &  42 \\
Near paragraph        & -- & --  &  7 &  38 \\
Over-broad holding    & -- & --  &  8 &  23 \\
Swapped party/outcome & -- & --  &  5 &  33 \\
Altered quote         & 15 & 35  &  8 &  42 \\
Wrong code (IPC/BNS)  &  4 & 20  & 13 &  12 \\
\midrule
\textbf{Total}        & \textbf{82} & \textbf{204} & \textbf{89} & \textbf{343} \\
\bottomrule
\end{tabular}
\end{table}

% ============================================================
\section{The Complexity-Adaptive Verification Pipeline}
\label{sec:pipeline}

The pipeline separates retrieval planning (adaptive depth) from answer auditing (evidence
verification), implemented as a LangGraph state machine. Baseline and proposed variants share the
same vector store (ChromaDB), embedding model (\texttt{text-embedding-3-small}), chat model
(\texttt{gpt-4o-mini}), and generation prompt, isolating each component's effect.

\textbf{Stage 1 --- Hybrid complexity scoring.}
\begin{equation}
S = \frac{0.35\,n_{\!needs} + 0.25\,r_{type} + 0.15\,m_{src} + 0.15\,\ell_{lex} + 0.10\,\lambda_{llm}}{\sum w_i}
\label{eq:complexity}
\end{equation}
$n_{needs}\!=\!\min(N\!-\!1,3)/3$ counts distinct information needs (LLM call, temperature 0);
$r_{type}$ maps reasoning type to $\{0, 0.5, 0.6, 0.7, 1.0\}$; $m_{src}\!\in\!\{0,1\}$ flags
multi-source queries; $\lambda_{llm}\!\in\!\{0, 0.5, 1.0\}$ is the LLM's own complexity label;
and
\begin{align}
\ell_{lex} &= 0.35\cdot\mathbf{1}[\text{cmp}] + 0.25\cdot\min\!\left(\tfrac{\text{enum}}{2},1\right) \notag\\
           &\quad + 0.25\cdot\min\!\left(\tfrac{\text{nest}}{3},1\right) + 0.15\cdot\min\!\left(\tfrac{|\text{ent}|-1}{3},1\right)
\label{eq:lexical}
\end{align}
is a purely deterministic lexical feature (no LLM call). $S$ is mapped to retrieval depth $k$ by
continuous interpolation: simple $k\!\in\![1,3]$, moderate $k\!\in\![4,7]$, complex $k\!\in\![8,12]$.

\textbf{Stage 2 --- Sub-question decomposition and RRF.}
When $N\!\geq\!2$, the system retrieves separately for the original query and up to four
sub-questions, then merges with Reciprocal Rank Fusion:
\begin{equation}
  \text{RRF}(d) = \textstyle\sum_{q \in Q} \frac{1}{60 + \text{rank}(d,\,q)}
  \label{eq:rrf}
\end{equation}
A coverage guarantee ensures the top chunk of each sub-query is included before filling by RRF
order, preventing bridge-hop chunks from being crowded out.

\textbf{Stage 3 --- Confidence and bounded expansion.}
\begin{equation}
  C = 0.5\,\tilde{s}_1 + 0.2\,\overline{\tilde{s}} + 0.3\,\text{cov}
  \label{eq:confidence}
\end{equation}
where $\tilde{s}_i$ normalises cosine scores and $\text{cov}$ is the fraction of sub-queries
covered. If $C < 0.5$ and budget remains, $k$ expands by 3 (capped at 15) for uncovered
sub-queries only.

\textbf{Stage 4 --- Evidence verification.}
Answers are split into atomic claims by an LLM call. Each claim is matched to 3-sentence evidence
windows by cosine similarity, then adjudicated by \texttt{cross-encoder/nli-roberta-base}:
\begin{align}
&\text{entail}(w,c)\geq 0.75 \Rightarrow \textsc{supported} \notag \\
&\text{contra}(w,c)\geq 0.75 \wedge \text{entail}<0.40 \Rightarrow \textsc{unsupported} \notag \\
&\text{otherwise} \Rightarrow \textsc{ambiguous} \to \text{LLM judge}
\label{eq:nli}
\end{align}
Ambiguous claims escalate to an LLM judge with a quote-check guard ($\geq$85\% fuzzy match).
Unsupported claims trigger bounded re-retrieval (max 2 retries); if still unsupported, an
explicit uncertainty notice replaces the claim.

% ============================================================
\section{Results}
\label{sec:experiments}

\textbf{Setup.} Items split 30/70 dev/test by source judgment (seed 13). B2 threshold tuned on
dev only (0.9 on v0; 0.7 on v1). Metrics: precision, recall, F1, false-flag (FF) rate on
faithful controls---all with 95\% bootstrap CIs (2,000 resamples). For MuSiQue: 36 questions
(12 per group) from the v1.0 dev set; corpus = 954 paragraphs; correctness judged by GPT-4.1
against gold; comparisons use paired bootstrap and McNemar tests (10,000 resamples).

\subsection{Legal Benchmark: Overall Detection}

Table~\ref{tab:overall} shows overall detection results. B0 (existence check) catches only
fabricated citations (F1 = 0.36 on v0, 0.31 on v1) and misses everything else. B1 (LLM-only,
\texttt{gpt-4o-mini}) reached F1 = 0.99 on v0---confirming v0 alone is too easy. B2 (local NLI)
dropped from F1 = 0.96 on v0 to 0.86 on v1, reflecting the additional difficulty of near-miss
types in v1 rather than a controlled regression.

\begin{table}[t]
\centering
\caption{Overall detection on the test split. 95\% bootstrap CIs (2,000 resamples).
         B1 evaluated on v0 only.}
\label{tab:overall}
\renewcommand{\arraystretch}{1.1}
\footnotesize
\begin{tabular}{@{}lllll@{}}
\toprule
\textbf{System} & \textbf{Prec.} & \textbf{Recall} & \textbf{F1} & \textbf{FF\%} \\
\midrule
\multicolumn{5}{l}{\textit{v0 (n=204)}} \\
B0 Existence  & 1.00 & 0.22 [.15,.29] & 0.36 [.27,.45] & 0.00 \\
B1 LLM        & 1.00 & 0.98 [.96,1.0] & 0.99 [.98,1.0] & 0.00 \\
B2 Local NLI  & 0.93 [.90,.97] & 0.99 [.97,1.0] & 0.96 [.94,.98] & 25\% \\
R  Date rule  & 1.00 & 0.12 [.08,.18] & 0.22 [.14,.30] & 0.00 \\
\midrule
\multicolumn{5}{l}{\textit{v1 (n=343)}} \\
B0 Existence  & 1.00 & 0.18 [.13,.23] & 0.31 [.24,.38] & 0.00 \\
B2 Local NLI  & 0.86 [.81,.90] & 0.87 [.82,.91] & 0.86 [.83,.89] & 31\% \\
R  Date rule  & 1.00 & 0.05 [.03,.08] & 0.10 [.05,.15] & 0.00 \\
\bottomrule
\end{tabular}
\end{table}

\subsection{Per-Type Analysis and Abstention}

Table~\ref{tab:pertype} shows per-type recall on v1. B2 handles blunt errors well (fabricated
case: 1.00; wrong paragraph: 0.98; wrong code: 1.00) but drops for near-miss types---especially
over-broad holdings (0.52 [0.30, 0.73]). This is expected: NLI is trained to accept general
claims entailed by their strict counterparts, but in law, omitting a qualifying condition is an
error, not a valid paraphrase.

The NLI baseline falsely flags 31\% of faithful controls (34/111). \textbf{Abstention helps.}
Reporting only on the most confident 20\% of items achieves 0.99 accuracy (AURC = 0.083),
falling to 0.93 at 40\% and 0.81 at full coverage.

\begin{table}[t]
\centering
\caption{Per-type recall on v1 (test split). 95\% bootstrap CIs.}
\label{tab:pertype}
\renewcommand{\arraystretch}{1.1}
\footnotesize
\begin{tabular}{@{}lrll@{}}
\toprule
\textbf{Type} & \textbf{n} & \textbf{B0} & \textbf{B2 Local NLI} \\
\midrule
Fabricated case       & 42 & 1.00 & 1.00 [1.00, 1.00] \\
Wrong paragraph       & 42 & 0.00 & 0.98 [0.92, 1.00] \\
Near paragraph        & 38 & 0.00 & 0.92 [0.83, 1.00] \\
Altered quote         & 42 & 0.00 & 0.81 [0.69, 0.92] \\
Swapped party/outcome & 33 & 0.00 & 0.76 [0.60, 0.91] \\
Over-broad holding    & 23 & 0.00 & 0.52 [0.30, 0.73] \\
Wrong code (IPC/BNS)  & 12 & 0.00 & 1.00 [1.00, 1.00] \\
\bottomrule
\end{tabular}
\end{table}

\subsection{MuSiQue: Adaptive Pipeline (Negative Result)}
\label{sec:negative}

Table~\ref{tab:musique} reports the ablation. The adaptive pipeline sets retrieval depth sensibly:
mean chunks are 1.9 / 5.1 / 8.2 for simple, moderate, and complex questions vs.\ a flat 5.0 from
the baseline; classifier-hop agreement is 64\%. However, better depth did not give better answers.
Fixed top-5 scored 0.583 correctness; A+V scored 0.444 (McNemar $p = 0.0625$, 5 vs.\ 0
discordant pairs). No variant beat the baseline on any metric. The likely reason: 954 paragraphs
(${\approx}20$ per question) is too small to discriminate retrieval policies; verification overhead
remains constant regardless. A larger corpus is needed.

\begin{table}[t]
\centering
\caption{MuSiQue ablation --- 36 questions (12 simple / moderate / complex).
         Mean [95\% CI]. McNemar $p$ vs.\ baseline B.}
\label{tab:musique}
\renewcommand{\arraystretch}{1.15}
\footnotesize
\begin{tabular}{@{}lp{1.9cm}p{1.9cm}rr@{}}
\toprule
\textbf{Variant} & \textbf{Correct} & \textbf{Recall@k} & \textbf{Lat.} & \textbf{McNem.} \\
\midrule
B~~Fixed top-5   & 0.583 [.417,.750] & 0.731 [.632,.829] & 1.6s  & ---  \\
A~~Adaptive      & 0.444 [.278,.611] & 0.662 [.549,.771] & 3.0s  & .0625\\
A$^-$ No decomp  & 0.444 [.278,.611] & 0.606 [.484,.722] & 2.8s  & .0625\\
V~~Verify only   & 0.500 [.333,.667] & 0.745 [.644,.843] & 10.4s & .2500\\
A+V Both         & 0.444 [.278,.611] & 0.699 [.588,.801] & 21.9s & .0625\\
\bottomrule
\end{tabular}
\end{table}

\subsection{Comparison with the Base Paper (Verma 2026)}

Table~\ref{tab:comparison} summarises how this work differs from Verma~\cite{verma2026citation}.
The key gaps we fill: (a) Indian jurisdiction with its unique IPC/BNS transition error; (b) seven
error types including three near-miss categories absent from prior work; (c) a local NLI verifier
requiring no API access; and (d) abstention analysis (AURC = 0.083). No direct numerical
comparison is made---the corpora, verifiers, and jurisdictions differ too much.

\begin{table}[t]
\centering
\caption{Comparison with Verma~\cite{verma2026citation} (base paper).}
\label{tab:comparison}
\renewcommand{\arraystretch}{1.1}
\footnotesize
\begin{tabular}{@{}p{1.3cm}p{2.7cm}p{2.7cm}@{}}
\toprule
\textbf{Aspect} & \textbf{Verma (2026)} & \textbf{This work} \\
\midrule
Corpus       & US opinions \& briefs     & 199 India SC judgments \\
Error types  & 2 (wrong case, pinpoint)  & 7; incl.\ 3 near-miss; IPC/BNS \\
Verifiers    & 14 LLM configs            & Local NLI, exist.\ check, date rule \\
Abstention   & Not addressed             & Risk-coverage; AURC = 0.083 \\
Best F1      & ${\sim}0.90$              & 0.99 (v0, LLM); 0.86 (v1, NLI) \\
\bottomrule
\end{tabular}
\end{table}

% ============================================================
\section{Discussion, Limitations, and Conclusion}
\label{sec:discussion}

\textbf{Why NLI fails on over-broad holdings.} NLI models trained on SNLI/MNLI assign high
entailment when a strict rule (``X only if Z'') appears to imply its generalisation (``X''). In
law, however, omitting the condition is an error, not a paraphrase. This is a fundamental
mismatch between legal citation support and the NLI training objective, not an implementation bug.
Legal-specific entailment objectives or ruling-aware classification heads are needed.

\textbf{Practical takeaways.}
(1)~Existence checks miss at least 69\% of citation errors---they are not sufficient for any real
legal AI deployment.
(2)~Confident abstention (top 20\%) achieves 0.99 accuracy with a fully local NLI model, making it
a practical strategy for high-stakes queries.
(3)~Adaptive retrieval only pays off on large, noisy corpora where retrieval is the actual
bottleneck. On small corpora, the depth-adaptation overhead dominates.
(4)~Over-broad holdings and swapped outcomes require capabilities beyond vanilla NLI---legal
semantic entailment and coreference tracking are open research problems.

\textbf{Limitations.} All items are LLM-generated; no lawyer has validated labels. We cover one
court, two years, 199 judgments, English only. The IPC/BNS mapping table is partial and unverified
against the official Gazette. MuSiQue uses only 36 questions on a small corpus with automatic
evaluation; both pipeline and judge are OpenAI-family and may share blind spots. Judgment texts are
not released; licence for redistribution is unconfirmed. This is a research prototype---not for
legal advice.

\textbf{Conclusion.} We presented \textbf{IndiaLegalCite}, a versioned citation-support benchmark
for Indian Supreme Court judgments with seven corruption types (three novel near-miss, one
IPC/BNS transition). Key findings: existence checks catch only fabricated citations; local NLI
reaches F1 = 0.86 on v1 but with high false-flag rate on over-broad holdings; abstention lifts
accuracy to 0.99 at 20\% coverage; and the adaptive RAG pipeline correctly modulates depth but
provides no accuracy gain at $6$--$14\times$ latency cost. Next steps: LLM checker on v1,
lawyer-validated labels, High Court data, and a larger MuSiQue corpus. Code and results at
\url{https://github.com/[repository]}.

% ============================================================
% REFERENCES
% ============================================================
\begin{thebibliography}{99}

\bibitem{magesh2025hallucination}
V.~Magesh et al.,
``Hallucination-free? Assessing the reliability of leading AI legal research tools,''
\textit{J. Empirical Legal Studies}, 2025; arXiv:2405.20362.

\bibitem{verma2026citation}
A.~Verma,
``Is this citation on point?''
arXiv:2608.12571, 2026.

\bibitem{qian2026gandr}
C.~Qian et al.,
``GANDR: Claim auditing for verifiable legal answer generation,''
arXiv:2609.10293, 2026.

\bibitem{jeong2024adaptiverag}
S.~Jeong et al.,
``Adaptive-RAG: Learning to adapt retrieval-augmented LLMs through question complexity,''
in \textit{Proc. NAACL}, 2024; arXiv:2403.14403.

\bibitem{asai2024selfrag}
A.~Asai et al.,
``Self-RAG: Learning to retrieve, generate, and critique through self-reflection,''
in \textit{Proc. ICLR}, 2024; arXiv:2310.11511.

\bibitem{yan2024crag}
S.-Q. Yan et al.,
``CRAG: Corrective retrieval augmented generation,''
arXiv:2401.15884, 2024.

\bibitem{min2023factscore}
S.~Min et al.,
``FActScore: Fine-grained atomic evaluation of factual precision in long form text generation,''
in \textit{Proc. EMNLP}, 2023; arXiv:2305.14251.

\bibitem{gao2023alce}
T.~Gao et al.,
``Enabling large language models to generate text with citations,''
in \textit{Proc. EMNLP}, 2023; arXiv:2305.14627.

\bibitem{gao2023rarr}
T.~Gao, A.~Fabbri, and C.~Manning,
``RARR: Researching and revising what language models say, using language models,''
in \textit{Proc. ACL}, 2023; arXiv:2210.08726.

\bibitem{pipitone2024legalbench}
N.~Pipitone and G.~H. Alami,
``LegalBench-RAG: A benchmark for retrieval-augmented generation in the legal domain,''
arXiv:2408.10343, 2024.

\bibitem{trivedi2022musique}
H.~Trivedi et al.,
``MuSiQue: Multihop questions via single-hop question composition,''
\textit{Trans. Assoc. Comput. Linguistics}, 2022.

\bibitem{indiansc2026}
``Indian Supreme Court Judgments,''
AWS Open Data registry (Dattam Labs). [Online]. Available:
\url{https://registry.opendata.aws/indian-supreme-court-judgments/}

\bibitem{paul2022inlegal}
S.~Paul et al.,
``Pre-training transformers on Indian legal text,''
in \textit{Proc. JURIX}, 2022; arXiv:2209.06049.

\bibitem{chalkidis2020legalbert}
I.~Chalkidis et al.,
``LEGAL-BERT: The muppets straight out of law school,''
in \textit{Findings of EMNLP}, 2020; arXiv:2010.02559.

\end{thebibliography}

\end{document}
